\documentclass{article}
\usepackage{graphicx} % Required for inserting images
\usepackage{mathtools}
\usepackage{indentfirst}
\usepackage{amsmath}
\usepackage{booktabs}   
\usepackage[backend=biber,notes,17th]{biblatex-chicago}
\addbibresource{referencesthesis.bib}
\usepackage{float}
\usepackage{placeins}

\title{ Pre-Thesis Assignment}
\author{Bruno Jose Fonseca}
\date{November 2025}

\begin{document}

\maketitle

\section{Proposal v1}
\subsection{Introduction}
I was inspired by Edward Miguel's Incentives to Learn Paper, where he ran an Randomized Control Trial based on a Girls Scholarship Program in two districts in western Kenya, Busia and Teso.\autocite[esp.~441]{kremer-miguel-thornton-2009} The merit scholarship program would grant the top 15th percentile best performing girls a 19\$ scholarship (KSh 1500), which is around 5\% of the GDP per capita of the time (400 dollars). \autocite[441]{kremer-miguel-thornton-2009} While strong and statistically significant positive results were observed in Busia, such were not seen in Teso. They stated that Teso's there is “a tradition of suspicion of outsiders in Teso, and this has at times led to misunderstandings with NGOs there.”\autocite[441]{kremer-miguel-thornton-2009}  The authors argued that events during the experiments in Teso affected the results, during the experiment there was a lightning strike that hit and destroyed a school not in the program killing 7 children, and “some community members associated the lightning strike with the NGO” as representatives were in the school related to another project just hours before the strike and argued that there were various religious connections to this phenomena. Right after, 6 schools in Teso pulled from the program, and there was even a student who turned down the scholarship after winning. 

Ultimately, this led me to question: To what extent does community collaboration influence the performance outcomes of international aid projects? To answer this question, I want to utilize the lightning strike as a way to measure the receptivity of the program. This is based on the assumption that schools near the lightning strike will be less willing to collaborate with the program than those further away. This is especially probable as this is an agricultural area with low use of cellphones and other means of communication other than word of mouth. Hence, the way the message was transmitted was probably incredibly dependent on distance, and also, the further away from the crash, the less the people would see the ruins of the schools and be impacted by the messaging and more likely they would be receptive. I would measure cooperation based on the proportion of students in the school who took the scholarship exam, as in Teso, “significantly more high-achieving students took the 2001 exam in comparison schools relative to program schools,” which suggests a boycott of the program. I would measure proportion as an expression of distance $P_i = \alpha + \delta\, d_i + \varepsilon_i$. If this is well supported, then I will estimate the impact on educational outcomes  $E_i= \alpha + \tau T_i+ \delta\,d_i +  \gamma (d_i*T_i) +\epsilon i$. 

Data is available and accessible. I have contacted Edward Miguel, and he believes it would be an interesting and worthwhile pursuit, and he has agreed to send me the data. The one thing I was worried about was how to measure the distance between schools; however, each school has GPS data associated with it, which I can use to accurately measure the geospatial variable. I am unsure whether the GPS location for the lightning strike is available; however, the authors included geographic visualizations in their article, suggesting that this information is available, can be extracted or reliably estimated from the maps if necessary. This would be sufficient for this investigation. 

Currently, access to the dataset is pending completion of the Social and Behavioral Research Investigators course, a prerequisite for handling the non-anonymized version of the data. Once this certification is completed, full access will be granted. The data have been used in a peer-reviewed article, hence the credibility and quality of the data are verified. 


\subsection{Literature Review}
This study would both contribute to a larger literature in the analysis of cooperation on economic and social outcomes, as well as contribute to quantifying and filling a gap directly related to the Incentives to Learn article. Most researchers have attempted to measure the impacts of implementing more collaborative designs when implementing development projects, such as the implementation of participatory committees and inclusion of community members (Mansuri and Rao 2013; Casey, Glennerster and Miguel 2012; White, Menon and Waddington 2018; Casey 2018).\autocites{mansuri-rao-2013}{casey-glennerster-miguel-2012}{white-menon-waddington-2018}{casey-2018-annrev} This is incredibly fascinating and is interesting as a means of project creation, and there have been positive results. Mansuri and Rao (2013) conducted a large meta-review of various World Bank projects and demonstrated that collaboration improves project outcomes; however, this translates weakly into welfare outcomes.\autocite[see esp.~ch.~2]{mansuri-rao-2013} They ultimately concluded that collaboration is complementary but not sufficient for an adequate project. White, Menon, and Waddington (2018) arrive at a similar conclusion through a mixed-method review focusing more on specific community-driven development, while the former defines and encompasses collaboration in a broader light, including the forming of government institutions. \autocite[]{white-menon-waddington-2018} 

Casey, Glennerster, and Miguel (2012) use the methodological robustness of an RCT with a pre-analysis plan to tackle this question, and it saw short-run positive gains in economic and educational outcomes.\autocite[]{casey-glennerster-miguel-2012} Although “collaboration” was delivered, it did not see long-term impacts and lasting changes in the organization and structure of the community. 
	
This demonstrates a faulty assumption made by most of these researchers, which is that even though the mechanisms for engagement might be there, the receptiveness for such mechanisms might not be well-received. This misses the point of receptiveness, as even with these interesting mechanisms, the project might not be well received. Which is an issue as the institutions and design are useless if not being explored in the way it was intended to do so. There is a difference between the collaboration structure and receptiveness, and the latter is necessary for the former, but in these cases, it was overlooked. 

This was addressed in Edward Miguel's Incentives to Learn paper, as he explicitly stated that socio-economic, religious, and natural factors impacted the educational outcomes and the receptiveness of the treatment, though this was not the main question he attempted to tackle.\autocite[]{kremer-miguel-thornton-2009} Furthermore, the impact of receptiveness was not quantified, and hence, the educational effects could not have been estimated in Teso in a scenario where there was adequate collaboration. The researchers did not test their own assumptions and arguments on why the educational outcomes were not successful in Teso, which leaves the question of the impact of receptiveness theorized, though not answered.

\subsection{My Approach}
So I identified issues of a lack of quantification and a lack of focus on the receptiveness variable. The development of this paper to the existing literature would be 3-fold. 

First, it would focus on receptiveness as a measure of collaboration and not structural project design. This, I argue, is a better approach than done beforehand, as even though there might be structures that encourage collaboration and participation by the community, it does not mean that the community is participating and receiving the project well. Furthermore, they are, to a certain extent, different measures in the policy project, which can be self-enhancing. Collaboration structures are established during the process where the project is being planned, while receptiveness can be enhanced in a project by choosing communities that are open to outside considerations or by conducting marketing programs to prepare communities for the project to come. Furthermore, if there is receptivity, then the collaborative structures in place will have a larger impact as they will be used by the community. 

I will measure this in a different way, while most of the cited studies use RCT or other methods comparing projects with collaborative structures with projects where they are not present, this design would not be possible to understand the variable of receptiveness. One cannot assign or not assign receptiveness, as that is, to a certain extent, something that is endogenous for the community; hence, using a natural experiment is a way to circumvent this issue. The lightning strike was completely random in time, randomly placed, and only had any impact in the study due to the reaction of individuals. Hence, the lightning, an exogenous variable, impacts the study only through the impact that it had on the receptiveness of the program. So by using distance to the lightning as a proxy for receptiveness then one can see the impact of it in the educational outcomes. 

Lastly, this dataset will bring a novel interpretation, as this is a natural experiment within an RCT. I will use the dataset from Miguel's Incentive to Learn, which has a comparison of treatment and control schools. I will measure the impact of the lightning both in the treatment and control schools, which will demonstrate if the lightning had any other impact on education outcomes other than through reluctance to participate in the scholarship program. Then I will specifically look at the effect of educational outcomes using an interaction between treatment and distance. Through this RCT and natural experiment approach, the randomized nature of the treatment will make it so that any other possible correlation with the lightning strike will be on average, be canceled out, and would affect control and treatment schools equally. This will make the effects of the lightning impact educational outcomes only through the responsiveness to the program and the willingness to engage with it. This use of a natural experiment within an RCT is innovative in this research, and it will be an important novel way to understand the impact of 

Such changes and situations would greatly enrich the literature as it would focus on the receptivity, a precursor for the collaboration design to be effective. Furthermore, it would provide a new methodological approach using a natural experiment to best measure the receptiveness variable, and lastly, it would be within a formal RCT where educational outcome measures are incredibly reliable. 

\section{Proposal v2}

\subsection{Description of Data}

I have currently gotten access to the data from Professor Miguel. His data included all the data from the Randomized Control Trial of the Girls Scholarship Program (GSP). The main database which was used for the paper has 11,728 observations of 32 variables. This main and final database is at the student level. Such 32 variables include indicator variables on if the school which the student is studying in is in Teso or Busia, educational exams for the RCT from 2001 to 2003 (normalized by district comparison group), baseline scores of the 2000 pre RCT, school id, as well as characteristics from the student on various socio economic aspects e.g. iron roof, notebooks, whether used notebooks. There is a variable, which is the proportion of students in student(i) school that speak ateso as their first language. Ateso is native to the Teso ethnic group, meaning that this could be an interesting measure of ethnic composition of each school. There is no data conducted during the RCT on the ethnicity of each student, just post RCT surveys. These surveys were not conducted in all individuals, but in representative samples. There is also one variable on whether the student is part of the restricted sample. 

There is however a lot more data in the GSP folder which is relevant. Particularly, GPS data at the school level. This is particularly important to construct the data on the distance from the lightning that each school is in km. I have constructed this database setting the school which was hit by the lightning (Korisai, school ID 916) as the reference point. This school was not in the RCT. I have since merged the distance from lightning of the school, based on the pupils' school id. This was then included into the main database commented above. This will be used in the analysis. 

The authors in the original study ran the regressions on a restricted sample, as not only did schools drop out after the lightning strike (3 treatment and 3 control schools) as well as there were schools with insufficient 2000, 2001 or 2002 scores. This reduces the number of schools to 116 with 7258 students (observations) with the same 32 variables being collected. I will probably use the restricted data, for similar reasons. Though the dropping of the 6 schools would be interesting to measure as they were done because of the lightning strike, it would be impossible to then connect to the test results, and hence it would not be useful to do so. For this reason they won't be considered in the econometric analysis. “The restricted sample contains data for 91\% of schools in the baseline sample. Students in program schools account for 51\% of the restricted sample.” \autocites{kremer-miguel-thornton-2009} 

There is another dataset that needs to be constructed for this thesis. I want to investigate the proportion of students in each individual school that took the exam. For that I would make an indicator variable on whether the exam is available (1 = yes 0 = no) and calculate the proportion in each school. This would provide the proportion of students who took the exam in each school. I would also transfer to this dataset the variable $ateso_percentage$. This is the school measure on the percentage of students who speak ateso (Teso's ethnic language) as their first language. This can be a proxy of ethnic groups. The lightning was interpreted by local religious leaders of the Teso ethnic group to be a sign for god, so this would be an important control. 

\subsection{Evaluation of Data}

This data seems sufficient, to tackle the questions at hand. There is a lot of data from a 3 year RCT that is able to be used in a variety of ways. This has data from 3 different years, meaning that I can check if the effects from the lightning change from year 1 of the RCT to year 2. 

	Though the data has limitations, there is no specific value of cooperation with the program and any possible measure would be an imperfect proxy to such. There is also significant data attrition, especially in the Teso district, as mentioned in the original paper. This is seen especially in the further years of the RCT as students move schools or choose not to take the exam. Furthermore though ateso might be an accurate proxy for ethnicity there is no information collected on the individual ethnicity of the student. There are two problems with this: first this is in the school level, its the percentage of students in the school, and then this value is given to every individual student in the school, so there is no specific data of the ethnicity of the student. Also its not a perfect proxy as there might be students who though not speak ateso as their first language, they still may be from the ethnic community and impacted by the anti NGO rhetoric. 
	
	Another limitation is that, as I am focusing in Teso, there is a significant decrease in observations from the original and even from the restricted when Busia is not considered. This reduces from 7258 observations to 2941 observations. Though even then, this data seems sufficient to achieve adequate conclusions, there might be complications, especially as considering possible missing values or other issues which may reduce the applicable observations even further. In some models I would be focusing on girls, as they are the main focus of the RCT which would reduce the number of observations significantly. Ultimately considering missing values there are some regressions which would be run with only ~800 or ~1000 observations, which though sufficient for a simple OLS, might not have sufficient power for an interaction term analysis, which I have researched needs to be between 4 to 16 times greater than what would be necessary for a simple analysis. \autocites{gelman-2018-need16}

\subsection{Econometric Models}
I want to investigate the proportion of students who don't take the 2001 exam in each school as a measure of collaboration. The model is: \[
P_i = \alpha + \lambda\, d_i + \beta t_i + \varepsilon_i
\]

Where $P_i$ is the proportion of students in school $i$ that has taken the 2001 exam, and $d_i$ is the distance in km from the lightning strike to that specific school. I am focusing on the 2001 exams, as the exams were collected in November 2001 and the lightning strike occurred in June 2001. This gives the time frame of 5 months for the lightning strike to have effect. It is probable that this time frame is sufficient as the 6 schools pulled out immediately after. Returning to the Econometrics, $\lambda$ would be the percentage point increase or decrease per km distance of the proportion of students in a school that took the 2001 exam. I intend to control for ateso, which is the proportion of students in $i$ school whose first language is ateso, the indigenous language of the Teso ethnic group. This is an important control, as the effect of the lightning strike is through the belief that it was a sign from god, or punishment, for the collaboration with the NGO or foreigners. This is a belief held and propagated through religious leaders of the Teso ethnic group. Hence if specific areas have larger concentrations of ethnic teso's, it is important that this is controlled for. This would answer the question as if you see a percentage point increase, in other words if $\lambda$ is positive and statistically significant, then that means that there was higher collaboration the further away from the lightning strike. This would demonstrate a correlation between the distance of the lightning and the collaboration of students. 

Another approach which focuses on the student level, and that can differentiate between participation in the treatment and control group is through an interaction term.
\[ A_{s\,t} = \alpha + c\,d_{s\,t} + y\,T_{s\,t} +  \lambda\ (d_{s\,t}* T_{s\,t})+\beta \,t_{s\,t}+h\, s_{s\,t-1}+ f\, M+ \epsilon_{s t}\]

In this case A is a dummy variable of whether a student has taken the 2001 exam, hence it will represent the percentage chance a given student would have taken the 2001 exam. There is an interaction term between distance (d) and treatment (T). Lambda will ultimately see what is the impact of distance specifically in the treatment group, and c would see the effect of the distance in the control group. Specifically they would predict what is the percentage point change per km on the chances of a given student taking the exam. Y would be, what is the percentage point change, on the probability that a given student takes the exam, based on their assignment to treatment and control group. In this case it is expected that cooperation would increase with greater distance in both treatment and control groups, however cooperation would greater in treatment groups with the distance, as they would be encouraged to participate given the Girls Scholarship Program. I am control for t, which is the percentage of students at a given school that have ateso as their first language, as a proxy for ethnic group of that region. I will further control for the base 2000 score, as individuals with higher base scores are more likely to take the 2001 exams. This is to make sure that this is controlled and that there are not specific schools in areas which have higher scores that can bias the main interest coefficient. This would demonstrate if there has been more cooperation the further away from the lightning strike, this can be used to see the effects of the distance on the project outcomes, to try to demonstrate a causal relationship. 

I want to also investigate who and where the people who did not take this exam are. In the study it states that “among high ability girls in Teso with a score of at least +1 standard deviation on the baseline 2000 exam, comparison school students were over 20 percentage points more likely to take the 2001 exam than program school students, and this difference is statistically significant at 95\% confidence over parts of this range.” I would argue that specifically these high achieving girls are closer to the lightning strike, as a sort of a boycott to the program. Here I would run a similar regression analysis, however now on the individual level with a dataset that focuses on such high achieving girls. Hence I would make a specific dataset only with high achieving girls, which score more then 1 SD above others. 

\[ A_{s\,t} = \alpha + c\,d_{s\,t} + y\,T_{s\,t} +  \lambda\ (d_{s\,t}* T_{s\,t})+\beta \,t_{s\,t}+h\, s_{s\,t-1}+ \epsilon_{s t}\]

This is exactly the same model as above, for the similar reasons which are established. The main coefficients are interpreted in the same way however instead of the percentage point change of a given student taking the test, it is the percentage point change of a high achieving girl to take the test. In this case it brings extra value for this issue as it provides insight into the mechanism which bring about the small effect in Teso. The authors argue that the highest achieving girls in Teso did not take the exams, in a way of a boycott. If this difference is seen mostly in the region surrounding the lightning then it can be concluded that the lightning had a large effect, and this boycott lead to lower educational and project outcomes which can also be seen. This is a regression I believe it would be a stretch due to sample size, by limiting to girls the and even more the high achieving girls, there is a heavy decrease from the original ~ 11,000. This might be an econometric model where not sufficient conclusions can be drawn due to the lack of sample size. I approximate that this sample would be about ~ 200 girls as by definition the above 1 sd is around 13\% of the original sample. This is very small for an interaction term. This can be mitigated by dropping the controls, and on average they should be the same between treatment and control groups, which would still make the coefficients relatively reliable. 

Lastly I would make a econometric model that would try to isolate the educational effects of the treatment from the lightning:

\[ E_{s\,t} = \alpha + \lambda\,d_{s\,t} + \beta\,T_{s\,t} +  \gamma\ (d_{s\,t}* T_{s\,t})+ \epsilon_{s t}\]

In this case the main coefficient of interest is gamma (y) the average difference in educational outcomes, specifically, standard deviation difference in test scores by district of control school between treatment and control group at a given distance to the lightning strike. If this is positive and statistically significant this would mean that further away from the lightning strike, you could see positive treatment effects, and that if it would have not been for the lightning strike, there could have been such positive effects in Teso as well. Lambda would be the average difference standard deviation difference between test scores normalized by district between treatment and control at km 0 from the lightning, and Beta would be the average standard deviation of the test score change per km further from the lightning of the control group. I am not including any particular control as through the nature of the RCT, on average treatment and control groups on any given distance from the lightning strike will have similar characteristics and the only difference would be the administration of the treatment. This, if positive and statistically significant, would demonstrate that indeed the more collaboration you have with a project, as the lightning could have only affected the control and treatment group through collaboration. I am afraid that as I keep decreasing and decreasing the sample, it might be smaller, especially considering that I am working with interaction terms. 

\section{Initial Results}
\FloatBarrier
\begin{table}
\centering
\begin{tabular}[t]{lrrrrrrr}
\toprule
Variable & N & Missing & Mean & SD & Min & Median & Max\\
\midrule
score\_2000 & 3952 & 3306 & 0.141 & 0.987 & -2.079 & 0.038 & 5.290\\
score\_2001 & 3748 & 3510 & 0.088 & 0.988 & -2.244 & -0.005 & 3.778\\
score\_2002 & 6320 & 938 & 0.079 & 1.011 & -3.887 & 0.001 & 4.699\\
missing\_score\_2001 & 7258 & 0 & 0.484 & 0.500 & 0.000 & 0.000 & 1.000\\
dist\_lig & 7258 & 0 & 23.193 & 10.841 & 1.453 & 23.129 & 42.886\\
\addlinespace
bs\_questionnaire\_result & 7258 & 0 & 0.776 & 0.417 & 0.000 & 1.000 & 1.000\\
treatment & 7258 & 0 & 0.497 & 0.500 & 0.000 & 0.000 & 1.000\\
teso & 7258 & 0 & 0.405 & 0.491 & 0.000 & 0.000 & 1.000\\
ateso\_percent & 7258 & 0 & 0.372 & 0.390 & 0.000 & 0.109 & 1.000\\
\bottomrule
\end{tabular}
\end{table}



\FloatBarrier
\printbibliography

\end{document}
